\documentclass[12pt]{article}
\usepackage[utf8]{inputenc}
\usepackage{amsmath}
\usepackage{amssymb}
\usepackage{graphicx}
\usepackage{booktabs}
\usepackage{listings}
\usepackage{xcolor}
\usepackage{hyperref}
\usepackage{geometry}
\geometry{margin=1in}

\lstset{
    language=Python,
    basicstyle=\ttfamily\small,
    keywordstyle=\color{blue}\bfseries,
    commentstyle=\color{green!60!black},
    stringstyle=\color{red},
    numbers=left,
    numberstyle=\tiny\color{gray},
    stepnumber=1,
    numbersep=5pt,
    backgroundcolor=\color{gray!10},
    frame=single,
    breaklines=true,
    breakatwhitespace=false,
    tabsize=4,
    showstringspaces=false
}

\title{Insurance Costs Calculation Documentation}
\author{}
\date{}

\begin{document}

\maketitle

\section{Introduction}

The \texttt{insurance\_costs.py} script calculates operational insurance costs for transmission line assets. Insurance premiums are based on insurable asset value (conductors, structures, converters) and are paid annually over the project lifetime.

\section{Method Overview}

\subsection{Purpose}

This script calculates insurance costs by:
\begin{enumerate}
    \item Determining insurable asset value from build costs (with contingencies)
    \item Applying premium rate based on construction type
    \item Calculating annual premium
    \item Computing lifetime nominal cost
    \item Calculating present value of annual premiums
\end{enumerate}

\subsection{Calculation Approach}

The insurance cost calculation follows this structure:

\begin{align}
\text{Insurable Value} &= \sum_{\text{insured components}} \text{Component Cost with Contingencies} \\
\text{Annual Premium} &= \text{Insurable Value} \times \text{Premium Rate} \\
\text{Lifetime Nominal Cost} &= \text{Annual Premium} \times \text{Project Lifetime} \\
\text{Present Value} &= \sum_{t=1}^{T} \frac{\text{Annual Premium}}{(1 + r)^t}
\end{align}

where $T$ is the project lifetime and $r$ is the real WACC discount rate.

\subsection{Inputs and Outputs}

\textbf{Inputs:}
\begin{itemize}
    \item Build costs (conductor, structure, converter) with contingencies
    \item Insurance configuration (premium rates, insurable components)
    \item Construction type (affects premium rate)
    \item Project lifetime
    \item Financing parameters (WACC for present value)
    \item Timeline (delay years, construction years for payment start)
\end{itemize}

\textbf{Outputs:}
\begin{itemize}
    \item Insurable asset value
    \item Annual premium
    \item Premium rate
    \item Lifetime nominal cost
    \item Present value of insurance costs
\end{itemize}

\subsection{Integration with Other Scripts}

This script integrates with:
\begin{itemize}
    \item \texttt{build\_costs.py} - Provides insurable asset values (see \texttt{build\_costs\_documentation.tex})
    \item \texttt{yaml\_loaders.py} - Loads project configuration and insurance parameters (see \texttt{utilities\_documentation.tex})
    \item \texttt{financial\_utils.py} - Financial calculations (see \texttt{utilities\_documentation.tex})
    \item \texttt{csv\_output\_manager.py} - Writes results to CSV (see \texttt{csv\_output\_manager\_documentation.tex})
\end{itemize}

\section{Key Functions}

\subsection{\texttt{calculate\_insurance\_costs(insurance\_yaml, conductor\_cost\_with\_contingencies, structure\_cost\_with\_contingencies, converter\_cost\_with\_contingencies, construction\_type, project\_lifetime)}}

Calculates operational insurance costs based on insurable asset value.

\textbf{Parameters:}
\begin{itemize}
    \item \texttt{insurance\_yaml} (dict): Loaded insurance YAML data
    \item \texttt{conductor\_cost\_with\_contingencies} (float): Conductor cost with contingencies
    \item \texttt{structure\_cost\_with\_contingencies} (float): Structure cost with contingencies
    \item \texttt{converter\_cost\_with\_contingencies} (float): Converter cost with contingencies
    \item \texttt{construction\_type} (str): Type of construction (overhead, underground, subsea)
    \item \texttt{project\_lifetime} (int): Project lifetime in years
\end{itemize}

\textbf{Returns:}
\begin{itemize}
    \item Dictionary containing:
    \begin{itemize}
        \item \texttt{insurable\_value} (float): Total insurable asset value
        \item \texttt{annual\_premium} (float): Annual insurance premium
        \item \texttt{premium\_rate} (float): Premium rate applied
        \item \texttt{nominal\_lifetime\_cost} (float): Total lifetime cost in nominal terms
    \end{itemize}
\end{itemize}

\textbf{Key Logic:}
\begin{lstlisting}
# Calculate insurable value based on which components are insured
insurable_value = 0.0
if components.get("conductors", True):
    insurable_value += conductor_cost_with_contingencies
if components.get("structures", True):
    insurable_value += structure_cost_with_contingencies
if components.get("converters", True):
    insurable_value += converter_cost_with_contingencies

# Get premium rate (construction type specific or default)
premium_rate = premium_by_type.get(
    construction_type.lower(), 
    ins_cfg.get("premium_rate", 0.002)
)

# Calculate annual premium
annual_premium = insurable_value * premium_rate
nominal_lifetime_cost = annual_premium * project_lifetime
\end{lstlisting}

\subsection{\texttt{main()}}

Main execution function that orchestrates the insurance cost calculation and output.

\textbf{Workflow:}
\begin{enumerate}
    \item Loads project technical details and constructs category identifier
    \item Loads build costs to get insurable asset values
    \item Loads insurance parameters
    \item Calls \texttt{calculate\_insurance\_costs()} to compute costs
    \item Calculates present value of annual premiums
    \item Displays results and writes to CSV
\end{enumerate}

\section{Example}

The following example demonstrates a typical usage of the insurance cost calculation:

\begin{lstlisting}
# Example: 500 MW Overhead AC line
# Category: "Overhead/AC/500MW/Advanced Aluminum/NA"
# 
# Inputs:
#   - Conductor cost (with contingencies): $55M
#   - Structure cost (with contingencies): $25M
#   - Converter cost: $0 (AC project)
#   - Premium rate: 0.2% (overhead construction)
#   - Project lifetime: 50 years
#   - Payment start: Year 3 (after 2-year construction)
#
# Calculation:
#   Insurable value = $55M + $25M + $0 = $80M
#   Annual premium = $80M * 0.002 = $160,000
#   Lifetime nominal = $160,000 * 50 = $8M
#   Present value = $160,000 * PV annuity (50 years, 3% discount) = $4.2M
#
# Output:
#   - Insurable value: $80M
#   - Annual premium: $160,000
#   - Lifetime nominal: $8M
#   - Present value: $4.2M
\end{lstlisting}

\textbf{Integration Example:}

\begin{lstlisting}
# When called from ctcc.py:
# 1. Script loads build costs from build_costs.py
# 2. Calculates insurance costs
# 3. Writes to CSV via CTCCOutputManager:
csv_manager = CTCCOutputManager()
csv_results = {
    "annual_premium": results["annual_premium"],
    "nominal_lifetime_cost": results["nominal_lifetime_cost"],
    "pv_total": insurance_pv,
    "insurable_value": results["insurable_value"],
    "premium_rate": results["premium_rate"],
}
csv_manager.add_insurance_costs(csv_results)
csv_manager.write_batch_summary()
\end{lstlisting}

\section{Key Implementation Details}

\subsection{Insurable Components}

The script allows configuration of which components are insurable (conductors, structures, converters). By default, all components are insured, but this can be customized in the YAML configuration.

\subsection{Premium Rate by Construction Type}

Different construction types may have different premium rates. The script first checks for a construction-type-specific rate, then falls back to a default rate if not specified.

\subsection{Payment Timing}

Insurance premiums start at Commercial Operation Date (COD), which is after the delay and construction periods. Premiums are paid annually for the project lifetime.

\subsection{AFUDC Treatment}

Operational insurance is not AFUDC-eligible because it is an operating expense, not a construction cost. Only construction-period costs are eligible for AFUDC capitalization.

\subsection{Dependency on Build Costs}

This script requires build costs to be calculated first (via \texttt{build\_costs.py}) to determine insurable asset values. It imports and calls \texttt{load\_costs()} from \texttt{build\_costs.py}.

\end{document}

